\section{Prerequisites and preliminaries on \texorpdfstring{\cstars{}s}{σC*-algebras}}\label{sec:preliminaries}
Throughout this paper, the proof techniques rely on machinery from category theory, particularly monads and symmetric monoidal categories (as well as monoidal monads).
We assume that the reader is familiar with basic category theory, including functors, natural transformations, limits and adjunctions. 
Prior exposure to monads is not required, though a general sense of them is helpful; for additional background on monads, readers may consult \cite[Sections~4.1 and~4.2]{borceux94handbook2}.
As for symmetric monoidal categories, the usual intuition of a ``tensor product'' should be sufficient.\footnote{Surveys on monoidal categories are abundant: \cite{selinger} is oriented towards applications, mostly in logic, while \cite{EGNO} is foundational for mathematicians interested in representation theory; \cite{milne_notes,borceux94handbook2,cwm} constitute good introductions that do not dwell too much on the finer details.}

This technology is adopted in the categories of (commutative) \cstar{}s and the subclass of (commutative) \cstars{}s. 

\begin{convention}\label{conv:unital}
Throughout this paper, unless otherwise specified (e.g., in the discussion around \autoref{lem:sigmaC_ideal}), \cstar{}s and \cstars{}s are always unital, and $*$-homomorphisms preserve the unit.
\end{convention} 

While the reader is not required to have a complete background on this theory, some basic knowledge of \cstar{}s may be used without further explanations: this includes, e.g., the notion of positivity and the fact that positive unital maps (and so also $*$-homomorphisms) are contractive, i.e.\ of norm $\le 1$~\cite[Proposition 2.1]{paulsen2002}.

Instead, no prior background on \cstars{}s is assumed.
The reader may, however, consult the books \cite{saito15monotone,pedersen2018automgroups}, as well as the article \cite{fritz2025categories}, where the choice of \cstars{}s is motivated from a probabilistic perspective.

\subsection*{\texorpdfstring{\cstars{}s}{σC*-algebras}}
	Recall that a \cstar{} carries a canonical partial order on its self-adjoint part, in which $a \le b$ means that $b-a$ is of the form $c^*c$.
A \newterm{\cstars{}} (also known as a \emph{monotone $\sigma$-complete \cstar{}}) is a \cstar{} $\A$ where every norm-bounded monotone increasing sequence $(a_n)$, i.e.\ such that $a_n \le a_{n+1}$ and $\norm{a_n}\le \lambda$, has a supremum $\sup_n^{\A} a_n$.\footnote{This requirement holds if and only if we restrict it to sequences of \emph{self-adjoint} elements. 
The only if direction is clear, while the if direction follows by noting that $a\le b$ implies that $a-b$ is self-adjoint, so that $a$ and $b$ have the same imaginary part: therefore, in a monotone increasing sequence all terms share the same imaginary part, so the general case reduces to the self-adjoint one upon subtracting it.
Below, the two conditions are used interchangeably.}
The superscript $\A$ in $\sup_n^{\A}$ is meant for clarity, and may be dropped when the context is sufficiently clear.

The slogan to keep in mind is that passing from \cstar{}s to \cstars{}s is the operator-algebraic counterpart of passing from topological to measurable structure: while commutative \cstar{}s correspond to compact Hausdorff spaces, commutative \cstars{}s correspond to \emph{abstract} measurable spaces, as we shall see at the end of this section. 
The intuition is that we allow a limiting procedure weaker than the uniform norm, analogous to how measurable functions are naturally closed under pointwise limits.

Between \cstars{}s, one generally considers \newterm{$\sigma$-normal} maps, i.e.\ those positive maps that respect this class of suprema.
Explicitly, a positive map $\phi\colon \A \to \B$ is $\sigma$-normal when, for every norm-bounded monotone increasing sequence $(a_n)$ in $\A$, the image sequence $(\phi(a_n))$ is again norm-bounded monotone increasing\footnote{This is immediate for positive maps between {unital} \cstar{}s, see \cite[Proposition 2.1]{paulsen2002}.} and
	\[
		\phi\Bigl(\sup\nolimits_n^{\A} a_n\Bigr) = \sup\nolimits_n^{\B} \phi(a_n).
	\]

The \starhom{}s (respectively, \cpu{}\footnote{This abbreviation stands for \emph{c}ompletely
\emph{p}ositive \emph{u}nital.} maps) that are $\sigma$-normal are called \newterm{\starshom{}s} (respectively, \newterm{\scpu{} maps}).
In this paper, we consider the categories as defined in \cref{tab:categories}.
\begin{table}[ht]
	\begin{tabular}{ccc}
		\textbf{Categories} & Objects    & Morphisms    \\
		\toprule
		$\Calg$             & \cstar{}s  & \starhom{}   \\\midrule
		$\Calgcpu$          & \cstar{}s  & \cpu{} maps  \\\midrule
		$\SigmaC$           & \cstars{}s & \starshom{}s \\\midrule
		$\SigmaCcpu$        & \cstars{}s & \scpu{} maps\\ \bottomrule
	\end{tabular}
	\caption{The categories considered in this paper (recall \cref{conv:unital}). The reader should imagine a parallel table where we add $\comm$ in front of each category name, in order to restrict to commutative algebras. For example, $\cCalg$ is the category of commutative \cstar{}s with \starhom{}s, and $\cSigmaCcpu$ is the category of commutative \cstars{}s with \scpu{} maps.}\label{tab:categories}
\end{table}

A \newterm{$\sigma$-subspace} in a \cstars{} $\A$ is a $*$-subspace $V$ having all suprema of norm-bounded monotone increasing sequences: $(a_n)\subseteq V$ implies $\sup_n^{\A}a_n \in V$. %\footnote{A $*$-subspace is simply a vector space closed under involution: $v\in V$ iff $v^*\in V$.}
Similarly, \newterm{$\sigma$-ideals} are those norm-closed two-sided ideals which happen to be $\sigma$-subspaces, and the \newterm{$\sigma$-closure} of a $*$-subspace $W$ is the smallest $\sigma$-subspace containing $W$.

For the technical developments of \cref{sec:tensor,sec:sigmaCP}, we also introduce \newterm{order limits}.
Briefly, one says that a sequence $(a_n)$ is \newterm{order convergent} in a \cstars{} $\A$ if there are four norm-bounded monotone increasing sequences of self-adjoint elements $(b_{n,k})_n$ for $k\in\lbrace 1,2,3,4\rbrace$ such that $a_n=(b_{n,1} - b_{n,2})+i(b_{n,3}-b_{n,4})$, where $i$ is the imaginary unit.
In this case, the order limit is given by
\[\textstyle
	\olim_n^{\A} a_n \coloneqq \sup_n^{\A} b_{n,1} - \sup_n^{\A} b_{n,2} + i \left(\sup_n^{\A} b_{n,3} - \sup_n^{\A} b_{n,4}\right).
\]
Several crucial properties of these limits are contained in \cite[Section 2.1]{saito15monotone}.

In particular, they are:
\begin{itemize}
	\item well-defined, so that $\olim$ of a sequence does not depend on the choice of the $b_{n,k}$'s: whenever we consider norm-bounded monotone increasing sequences of self-adjoint elements $(b_{n,k})_n$ and $(b_{n,k}')_n$ for $k=\lbrace 1,2,3,4\rbrace$ such that
	      \[
		      a_n=(b_{n,1} - b_{n,2})+i(b_{n,3}-b_{n,4})=(b_{n,1}' - b_{n,2}')+i(b_{n,3}'-b_{n,4}') ,
	      \]
	      then
	      \begin{multline}
								\textstyle
		      \sup_n^{\A} b_{n,1} - \sup_n^{\A} b_{n,2} + i \left(\sup_n^{\A} b_{n,3} - \sup_n^{\A} b_{n,4}\right) =\\ \sup_n^{\A} b_{n,1}' - \sup_n^{\A} b_{n,2}' + i \left(\sup_n^{\A} b_{n,3}' - \sup_n^{\A} b_{n,4}' \right).
							\end{multline}
	\item linear: $\olim^{\A}_n (a_n + a_n')=\olim^{\A}_n (a_n)+ \olim^{\A}_n (a_n')$ and $\olim^{\A}_n(\lambda a_n) = \lambda \olim^{\A}_n a_n$ for all $\lambda \in \bbC$;
	\item $*$-preserving: $\olim^{\A}_n (a_n^*)=\olim^{\A}_n(a_n)^*$;
	\item compatible with multiplication: $\olim^{\A}_n (z a_n w) = z(\olim^{\A}_n a_n)w$ for all $z,w \in \A$.
\end{itemize}
Note that suprema do not respect these properties: for example, $-\sup^{\A}_n a_n $ becomes an infimum $\inf^{\A}_n (-a_n)$.
More instructively, multiplying two monotone increasing sequences may not yield a monotone increasing sequence because already the product of two positive elements is not necessarily positive.
Instead, 
\[ \olim^{\A}_n(a_n \olim^{\A}_m b_m)= \olim^{\A}_m (\olim^{\A}_n a_n b_m)=\olim^{\A}_n (a_n)\olim^{\A}_m(b_m).\]

It is a simple check that \scpu{} maps are exactly those cpu maps that respect order limits.
In contrast, it is currently unclear whether a $\sigma$-subspace is necessarily closed under order limits (see \cite[Section 3.5]{fritz2025categories}).

\subsection*{\texorpdfstring{\bsigma{}s}{Boolean σ-algebras}}
Another crucial aspect of \cstars{}s is their relation to \emph{abstract} measurable spaces, better known as \bsigma{}s.
Briefly, these are Boolean algebras with countable joins and meets, and for this reason they mimic $\sigma$-algebras.
The interested reader may refer to \cite{sikorski69boolean} for additional details.
Nonetheless, this paper does not employ the technology of Boolean algebras, and references will be provided when needed.
The material below is needed primarily in \cref{sec:giry,sec:sigmaCP} and may be skipped otherwise.

\begin{notation}
	We denote Boolean $\sigma$-algebras with $\Sigma$ and their elements with $U,V,U_n, \dots$, to remind the reader of the case of measurable spaces.
	Nonetheless, we use the Boolean algebra notations for operations: $\vee$ denotes the join, $\wedge$ the meet, $\neg$ the negation, $\bot$ the bottom element and $\top$ the top.
\end{notation}

It was shown in \cite[Theorem 4.1.34]{fritz2025categories} that $\cSigmaC$ is equivalent to the category of Boolean $\sigma$-algebras and $\sigma$-homomorphisms (here meaning Boolean algebra homomorphisms respecting countable joins and meets).
The \cstars{} associated to the \bsigma{} $\Sigma$ is denoted by $\Linfbool(\Sigma)$.
The notation of $\Linfbool(\Sigma)$ (``abstract $\Linf$'') is chosen to suggest that $\Linfbool(\Sigma)$ is an abstraction of the `concrete' \cstars{} associated to a measurable space $(X,\Sigma_X)$, given by
\[
	\Linf(X)\coloneqq \lbrace f\colon X \to \bbC \text{ bounded measurable}\rbrace
\]
and equipped with pointwise operations.
In fact, $\Linfbool(\Sigma_X)\cong \Linf(X)$ (\cite[Lemma 4.1.14]{fritz2025categories}).
More explicitly, the elements of $\Linfbool(\Sigma)$ are encoded by analogy with preimage maps of bounded measurable functions:
\[
	\Linfbool(\Sigma)\coloneqq \lbrace f\colon B(\bbC)\to \Sigma \mid f(U)=\top \text{ for some bounded }U \in B(\bbC) \rbrace,
\]
where $B(\bbC)$ denotes the Borel $\sigma$-algebra of $\bbC$.
From a commutative \cstars{} $\A$, one can show that $\A \cong \Linfbool(\Sigma)$, with $\Sigma$ its set of projections (which turns out to be a \bsigma{}).
The reader may refer to \cite[Section 4.1]{fritz2025categories} for further details.

From a probabilistic perspective, the category $\Stoch$ of measurable spaces and Markov kernels embeds fully faithfully (contravariantly) into  $\cSigmaCcpu$ (see \cite[Proposition 4.3.6]{fritz2025categories}). Accordingly, $\sigma$-cpu maps $\Linfbool(\Sigma)\to \Linfbool(\Sigma')$ may be viewed as \emph{abstract Markov kernels}.
While tangential to the main discussion, this perspective is included for readers interested in the connection with classical probability.

